\section{A section on the entire reduced boundary}
\label{floor:sec}

This section proves the adjunction statement needed on the reduced-boundary
model.  The distinction between a section on one component and a section on
the whole reduced boundary is essential.  In particular, none of the gluing
below follows merely from abundance on the normal components.

\begin{proposition}[Whole-boundary nonvanishing]
\label{floor:section}
Let $(V,D)$ be the compact K\"ahler dlt fourfold supplied by
Proposition~\ref{bir:reduced-model}, where $V$ is globally strongly
$\Q$-factorial and $D\ne0$ is reduced.  Put
\[
 J=K_V+D.
\]
If the actual rational adjoint line bundle $J$ is analytically nef, then,
for some positive integer $m$ for which $mJ$ is Cartier on $V$,
\begin{equation}
 \label{floor:goal}
 H^0\bigl(D,\OO_D(mJ)\bigr)\ne0.
\end{equation}
\end{proposition}

We prove the proposition by constructing sections on the normalized
components with matching canonical residues.  All adjunction boundaries
in the argument are the full effective rational differents; their
fractional parts are never discarded.  All sufficiently divisible degrees
will also be even.  Evenness removes the sign obtained by interchanging
two residue operations, but does not permit a change of the actual
adjunction line bundle.

\subsection{Adjunction and the descent locus}

Write $D=\sum_iD_i$ and let $Y_i=D_i^\nu$ denote the normalization of a
component.  Dlt chain adjunction gives an effective rational dlt pair
$(Y_i,\Gamma_i)$ and an equality of actual rational line bundles
\begin{equation}
 \label{floor:component-adjunction}
 J_i:=J|_{Y_i}=K_{Y_i}+\Gamma_i.
\end{equation}
Here and below restriction includes pullback by the indicated
normalization.  Further adjunction is always performed with the full
boundary in \eqref{floor:component-adjunction}.

We use the good SNC opens and chain-adjunction construction of
Lemma~\ref{bir:stratum-adjunction}.  Every lc center is
generically a stratum of distinct coefficient-one primes.  In the
particular modification used here, this can also be seen by following a
discrepancy-zero place from the starting SNC resolution: strict
discrepancy increase over a nonisomorphism locus, as in
Lemma~\ref{bir:stratum-comparison}, forces the general point
of its center to remain unchanged at every step.  The distinct primes at
that stratum consequently remain distinct.

Keeping only one component $D_i$ in the ambient boundary gives a plt
pair.  Indeed, any exceptional place with zero log discrepancy for
$(V,D_i)$ would also have zero log discrepancy for $(V,D)$ and would
have its center generically in one of these unchanged SNC strata.  For a
single smooth coefficient-one branch there is no such exceptional
zero-discrepancy place.  Comparisons are legitimate because the
components are effective rationally Cartier divisors.  Adjunction for
$(V,D_i)$ therefore supplies an effective rational klt pair on $Y_i$.
In particular, $Y_i$ has rational singularities.  A Moishezon $Y_i$ is
projective by Namikawa's theorem, since it is also compact K\"ahler
\cite{NamikawaProjectivity}.

\begin{lemma}[Codimension-one matching suffices]
\label{floor:descent}
The reduced space $D$ is $S_2$ and has only smooth points and ordinary
double crossings in codimension one.  Its codimension-one conductor
branches are the normalizations of the primes of
$\lfloor\Gamma_i\rfloor$.  They are paired between distinct components
$D_i,D_j$.  For a sufficiently divisible even $m$, a tuple
\[
 s_i\in H^0(Y_i,\OO_{Y_i}(mJ_i))
\]
descends to $H^0(D,\OO_D(mJ))$ if its restrictions agree on every paired
normalized conductor surface under the canonical adjunction
identifications.
\end{lemma}

\begin{proof}
First, both $\OO_V$ and $\OO_V(-D)$ are Cohen--Macaulay.  The first
assertion follows from klt rationality: decreasing the reduced dlt
boundary makes the zero-boundary space klt.  For the second, work
locally and trivialize a Cartier multiple of the integral Weil divisor
$D$.  The resulting normal cyclic index cover is \emph{\'etale in
codimension one}: divisorial orders in the defining equation are
multiples of the index.  The cover is klt by the finite canonical
pullback and discrepancy formula, hence Cohen--Macaulay.  The desired
divisorial sheaf is an eigensummand of the pushforward of its structure
sheaf, so is Cohen--Macaulay as well.  Normality identifies the reduced
ideal of $D$ with $\OO_V(-D)$.  The sequence
\[
 0\longrightarrow\OO_V(-D)\longrightarrow\OO_V
   \longrightarrow\OO_D\longrightarrow0
\]
now shows that $D$ is Cohen--Macaulay, in particular $S_2$.

For the codimension-one description, take general transverse surface
germs in the dlt adjunction calculation.  At a single coefficient-one
branch the plt local model is a cyclic quotient of a smooth pair, and
its boundary curve is smooth.  At two coefficient-one branches, the
different and the sub-lc condition give an ordinary double SNC point;
these are precisely the generic deeper lc strata described above.  The
branches come from distinct primes.  To lift the assertion about a
single-branch transverse curve to regularity at the corresponding
point of $D$, note that the transverse boundary curve is generically
reduced and has no embedded points by Cohen--Macaulayness.  Its
regularity lifts by lifting generators of the maximal ideal through
the transverse parameters.  This proves the asserted codimension-one
description.

Let $S$ be a normalized conductor surface on $Y_i$.  Chain adjunction
gives
\begin{equation}
 \label{floor:surface-adjunction}
 J_S=K_S+\Xi_S=J|_S.
\end{equation}
The matching branch on $Y_j$ is the normalization of the same image,
and the two normalizations agree.  At a general SNC point the two
iterated residue maps differ by the sign from interchanging two
differentials.  Their even powers are equal.  On a common smooth model
the two invertible pullback subsheaves of meromorphic pluricanonical
forms are therefore equal: they have the same invertible domain and
the same meromorphic map on a dense open.  Thus this is an equality of
the actual residue lines, not only a numerical comparison.  The same
argument applies to further normalized lc strata.  Proper subcenters
are generically deeper SNC strata, so the chains on either side
identify the same further centers and the same even residue maps.

At an ordinary double point the usual two-branch local calculation
glues sections precisely when their values agree on the intersection.
Hence a matching tuple is a section of the ambient invertible
restriction away from a subset of codimension at least two in $D$.
For a reduced $S_2$ analytic ring, regularity of a meromorphic element
is tested in its height-one localizations.  Applied after trivializing
$\OO_D(mJ)$, this extends the section across the omitted subset.
Equivalently, the class of the normalization-pushforward section
modulo the invertible sheaf has support of codimension at least two
and vanishes by the $S_2$ Hartogs criterion.  The residue comparisons
used here are the comparisons induced by the one ambient line bundle;
no arbitrary scalar choices in adjunction enter the descent.
\end{proof}

\begin{lemma}[Systems on the normal components]
\label{floor:component-systems}
Each $J_i$ is semiample as an actual rational line bundle.  There is a
holomorphic connected-fiber map to a normal projective variety and an
ample rational line bundle such that
\begin{equation}
 \label{floor:stein-system}
 f_i:Y_i\longrightarrow Z_i,\qquad J_i=f_i^*H_i.
\end{equation}
\end{lemma}

\begin{proof}
The normal compact K\"ahler threefold log-abundance theorem applies to
the effective rational adjunction pair
\eqref{floor:component-adjunction}: its actual adjoint is analytically
nef.  One may first pass to a projective crepant dlt modification and
use the $\Q$-factorial dlt form of the theorem.  We use here
Das--Ou's normal threefold theorem, together with their dlt-modification
theorem \cite[Corollary~1.3]{DasOuAbundanceII}
\cite[Theorem~5.2]{DasOuAbundanceI}.  Generation on the modification
descends by projection to the normal target.  These inputs concern
normal pairs; no nonnormal gluing assertion is being invoked.

For completeness, the passage to divisor representatives in these
theorems does not impose an additional global-frame hypothesis.  One
can first produce a section of a power of $J_i$ as follows.  If $Y_i$
is projective, use projective threefold log abundance.  Otherwise
resolve it by a smooth compact K\"ahler space.  If that space is
non-uniruled, smooth threefold canonical nonvanishing gives a section
which pushes to $J_i$, since $\Gamma_i\ge0$.  If it is uniruled and
non-Moishezon, take a resolved rational quotient.  Its smooth
non-uniruled compact K\"ahler base has nonnegative Kodaira dimension,
and its very general smooth quotient fibers are projective of
dimension at most two.  On a simultaneous resolution of the pair,
use the strict boundary and all reduced exceptional divisors as an
effective SNC boundary $\widetilde\Gamma$.  Log canonicity gives
\[
 K_{\widetilde Y}+\widetilde\Gamma
   =p^*J_i+E,\qquad E\ge0\quad\hbox{$p$-exceptional}.
\]
This adjoint restricts pseudo-effectively on a very general quotient
fiber.  Projective log nonvanishing in dimensions one and two supplies
fiberwise sections.  Assumption~\ref{hyp:campana}, including unit
coefficients and its invariant-base interpretation, now supplies a
section upstairs exactly as in the rational-quotient reduction.  It
pushes to a section of an actual power of $J_i$ by normality and
reflexivity.  The rational-quotient and uniruledness inputs are used in
their smooth compact K\"ahler scopes
\cite{CampanaRationalQuotient,OuUniruledness}.

Cancelling the boundary contribution in such a section gives a
meromorphic pluricanonical tensor.  Its divisor divided by its degree
is a rational canonical representative, compatible with pullback to
resolutions.  The discrepancies are the local canonical
discrepancies, and the divisible systems are still the actual
systems of $J_i$.  Thus the divisor formulation of threefold
abundance gives precisely the claimed semiampleness.  The Stein
factorization of a sufficiently high generated system yields
\eqref{floor:stein-system}.
\end{proof}

\subsection{Dominant clusters and a finite-product construction}

Set $r_i=\dim Z_i$ and $r=\max_i r_i$.  At a vertex with $r_i=r$,
call a normalized conductor surface $S$ \emph{dominant} when it
surjects onto $Z_i$.  By \eqref{floor:surface-adjunction} and
\eqref{floor:stein-system}, this is equivalent to
$\kappa(S,J_S)=r$.  On the opposite side of a matched conductor
surface, the same equality forces $r_j=r$, and the surface is
dominant there too.  Thus dominance is symmetric.

Form the finite graph whose vertices have $r_i=r$ and whose edges
are the dominant conductor surfaces.  Fix a connected component
of this graph, called a \emph{cluster}.  It is enough to construct a
nonzero matching tuple on this cluster which vanishes on every
nondominant branch, including nondominant edges within the cluster;
we then put zero on every other component.  The finitely many
nondominant images are proper subvarieties of the relevant $Z_i$.
An ample power has a nonzero section vanishing on their union.
An isolated vertex is therefore immediate.  If $r=0$, there are no
nonempty nondominant branch images.  If the cluster has an edge,
then $r\le2$, since an edge is a surface.  This also settles any
vertex with $r=3$.

We shall use normalized strata with the full chain adjunction
\[
 (L,\Delta_L),\qquad J_L=K_L+\Delta_L=J|_L.
\]
Each chosen object $L$ dominates the corresponding $Z_i$ and has
semiample $J_L$.  An arrow between objects means a $B$-bimeromorphic
comparison: on a common resolution, the actual adjunction lines and
their meromorphic pluricanonical pullback maps agree.  In sufficiently
divisible even degrees, arrows consequently give isomorphisms of
section spaces.  These isomorphisms commute with multiplication.

\begin{lemma}[Finite products of transported sections]
\label{floor:products}
Consider finitely many such objects and invertible arrows.  Suppose
that in one sufficiently divisible even degree $m$ the image of
every group of loops on its section space is finite.  At each object
choose a nonzero section in that degree, obtained by restricting a
base section which vanishes on the prescribed nondominant images.
Then, in a common higher degree, there are nonzero sections at all
objects, invariant under every arrow, and retaining the prescribed
vanishing.
\end{lemma}

\begin{proof}
Fix one orbit of objects.  At an object $L$, take the set of all
sections transported to $L$ from all chosen initial sections in that
orbit, along all paths.  The set is finite: there are finitely many
starting objects, and two paths from the same object differ by a
loop, whose image is finite.  Arrows biject these finite sets.
Multiply all their distinct members.  The resulting section is
nonzero because the object is irreducible, and the products are
carried to one another because pullback is multiplicative.  They
have equal degree within the orbit.  The initial section at $L$
occurs among its factors, so its prescribed zeros remain.  Taking
further powers makes the degrees equal across the finitely many
orbits.
\end{proof}

In each application with $r>0$ we shall verify that the invariant
sections on the objects descend to sections of the node's base
polarization.  This will also preserve the required vanishing: a
dominating object's section contains the pullback of the initial
vanishing base section as a factor.

We record the residue calculation used to construct internal arrows.
Suppose a log rational-curve fibration has, on its general fiber,
two distinct coefficient-one points.  The total horizontal weighted
degree of the boundary is two, so there are no other horizontal
markings.  After ordering the points, choose a rational coordinate
$x$ with them as zero and infinity.  A local base frame of an
adjoint power has, in relative log differentials, the form
\begin{equation}
 \label{floor:dlog-residues}
  a\,(d\log x)^{\otimes m},
\end{equation}
where $a$ is constant on the compact general rational fiber.  The
residues at zero and infinity differ by $(-1)^m$, and hence agree
for even $m$.  A change $x\mapsto c\,x$ adds a base differential,
which disappears in the top-form expression.  This proves that
pairing the two marks preserves restrictions of base sections.
For two degree-one branches it gives a bimeromorphic map between
their normalizations; for one degree-two branch it gives the
involution exchanging the two generic sheets.  The latter is
defined by normalization of the other main component of the
relative fiber square.  Since the lines on a common resolution
pull back the same base line, equality of the meromorphic residue
maps on this dense open proves the full $B$-crepant comparison.

\subsection{Clusters containing a non-Moishezon vertex}

We first describe the geometry of such a vertex; as before we omit
its subscript and write $(Y,\Gamma)$, $J_Y=f^*H$, and $Z$.

\begin{lemma}[The rational-quotient surface]
\label{floor:quotient-surface}
At a non-Moishezon vertex with a dominant conductor branch there is
a diagram of compact K\"ahler spaces
\begin{equation}
 \label{floor:quotient-diagram}
 \begin{tikzcd}[column sep=large]
 W \arrow[r,"g"] \arrow[d,"p"'] & P \arrow[d,"u"]\\
 Y \arrow[r,"f"'] & Z
 \end{tikzcd}
\end{equation}
in which $W$ is smooth and resolves the pair, $P$ is a smooth
nonprojective compact K\"ahler surface, both $g$ and $u$ have
connected fibers, the very general $g$-fiber is $\PP^1$, and every
$p$-exceptional divisor is vertical over $P$.  The horizontal
crepant boundary on $W$ is effective, with weighted degree two on
a general $g$-fiber.  Moreover $r\le1$ and
$H^0(P,K_P)\ne0$.
\end{lemma}

\begin{proof}
On an initial smooth projective log resolution of $Y$, restrict the
crepant equality
\[
 K_W+\Gamma_W=p^*J_Y
\]
to a very general smooth connected fiber of the map to $Z$.
Its right side is trivial.  Pair with the appropriate power of the
pullback of a K\"ahler class from $Y$.  Exceptional terms have zero
pairing, because their images have codimension at least two on this
general fiber, or they do not meet it.  Strict boundary terms have
nonnegative pairing; a dominant coefficient-one branch gives a
strictly positive term.  Hence the canonical class of this smooth
fiber is not pseudo-effective.  Ou's theorem implies that the
fiber is uniruled \cite{OuUniruledness}.  For $r=0$ this argument is
applied to the total space.

Take the almost holomorphic rational quotient of the smooth
K\"ahler resolution.  Its fiber has positive dimension, since the
rational curves through very general points of the fibers just
considered are quotient-contracted
\cite{CampanaRationalQuotient}.  Its base cannot have dimension
at most one.  Indeed, rational connectedness of the general quotient
fibers and the differential sequence would then force every
holomorphic two-form on a resolved total space to vanish.  The
K\"ahler projectivity criterion would make that total space
projective, a contradiction.  The quotient base is therefore a
surface, and the general quotient fiber is a smooth rational curve.
Those curves are vertical over $Z$, so $f$ factors meromorphically
through the quotient.  Resolving the graphs and the base gives
\eqref{floor:quotient-diagram}.  Stein factorization and very
general connectedness give connected $g$-fibers; the fibers of $u$
are connected as images of the connected fibers of $fp$.

If $P$ were projective, the rational-curve fibration would have
Moishezon total space.  For example, coherence and GAGA give
meromorphic sections of $g_*\OO_W(-K_{W/P})$ generating its general
fibers.  Evaluation embeds the general smooth rational fiber, and
these sections together with base functions give full algebraic
dimension.  Thus $P$ is nonprojective.

Every exceptional prime of the initial projective resolution of
$Y$ is projective over a compact curve or a point, hence Moishezon.
It cannot dominate $P$: a dominant generically finite map of
surfaces preserves algebraic dimension, by the norm or
characteristic-polynomial argument after finite Stein
factorization.  The strict transform of a Moishezon prime remains
Moishezon.  The further graph resolutions can be made isomorphisms
over a general quotient-base open, so their new exceptional primes
are vertical too.  Consequently all horizontal crepant boundary
terms are strict transforms of effective boundary components.
Adjunction on the general rational fiber makes their weighted
degrees sum to two.

A nonprojective compact K\"ahler surface has a nonzero holomorphic
two-form, again by the K\"ahler projectivity criterion.  Finally,
if $r=2$, then $P\to Z$ would be generically finite over a
projective surface, forcing $P$ to be Moishezon and projective.
Thus $r\le1$.
\end{proof}

We use two elementary consequences of nonprojectivity for surfaces.
If $Q$ is a smooth nonprojective compact K\"ahler surface, then the
intersection form on
\[
 N_Q=\NS(Q)_{\R}
\]
is nonpositive.  Indeed, a positive-square vector in this rational
subspace can be approximated by a rational one, with sign chosen to
have positive K\"ahler degree.  The corresponding line bundle has
quadratic section growth by Riemann--Roch: its high powers have no
top cohomology by Serre duality and the K\"ahler degree test.  This
would make $Q$ Moishezon and hence projective.  Moreover, if
$Q\to C$ is a map to a projective curve, $Q$ has no multisection.
For a curve $B$ dominating $C$, the class $B+tF$, with $F$ the
rational fiber class and $t\gg0$, would have positive square.

\begin{lemma}[Horizontality, including fractional boundaries]
\label{floor:horizontality}
Every dominant conductor branch at the non-Moishezon vertex of
Lemma~\ref{floor:quotient-surface} is horizontal over $P$.  A
cluster containing such a vertex consists entirely of
non-Moishezon vertices.  At a vertex there are at most two dominant
branches, counted with their generic degrees over $P$.
\end{lemma}

\begin{proof}
If $r=1$, a vertical surface which dominates $Z$ would have image
a curve in $P$ dominating $Z$.  This contradicts the preceding
no-multisection observation.  It remains to prove horizontality
when $r=0$, so that $J_Y\sim_\Q0$.

Choose a nonzero $\eta\in H^0(P,K_P)$.  On the general rational
fiber write the distinct horizontal marked points and their
coefficients as
\[
 b_1,\ldots,b_q\in\Q,\qquad
 0<b_j\le1,\qquad \sum_{j=1}^q b_j=2.
\]
Take a smooth compact K\"ahler generically finite base change
$b:P^+\to P$ splitting and ordering the markings.  It can be
constructed from a main component of a fiber product of the
horizontal prime normalizations over $P$, on the finite \emph{\'etale}
locus parametrizing an ordering of the distinct marks, and then
resolved.  Resolve the main component of the pulled-back rational
fibration, obtaining
\[
 h:W^+\longrightarrow W,\qquad
 g^+:W^+\longrightarrow P^+.
\]
Let $H_j$ be the prime closures of the ordered generic sections.
The very general rational fibers are unchanged by these resolutions.

For $j\ne k$ there is a meromorphic top form
\begin{equation}
 \label{floor:weighted-dlog}
 \beta_{jk}=d\log x_{jk}\wedge(g^+)^*b^*\eta,
\end{equation}
where $x_{jk}$ has horizontal divisor $H_j-H_k$.  This description
is valid meromorphically even across degenerations.  To see its
local construction on the base, the coherent sheaf
$(g^+)_*\OO_{W^+}(H_j-H_k)$ has generic rank one, because its
restriction to a general rational fiber has degree zero and is
trivial.  Divide the canonical meromorphic section by the
evaluation of a locally chosen generically nonzero direct-image
section.  Different choices change $x_{jk}$ by a base function on
the general fibers.  Its logarithmic differential wedges to zero
with the pulled-back top form of $P^+$, so
\eqref{floor:weighted-dlog} is independent of the choices.
Reversing $j,k$ changes its sign.

The rational vector $(b_j)$ lies in the convex hull of the vectors
$\mathbf e_j+\mathbf e_k$ for $j<k$.  These are exactly the
vertices of the rational polytope
\[
 \bigl\{(x_j):0\le x_j\le1,\ \sum_jx_j=2\bigr\}.
\]
Average a rational convex decomposition over the permutations
which preserve the coefficients.  After clearing denominators
and multiplying by an even integer, we obtain nonnegative even
integers $l_{jk}$ and an integer $m>0$ such that
\begin{equation}
 \label{floor:incidence-weights}
 \sum_{j<k}l_{jk}=m,\qquad
 \sum_{k\ne j}l_{\min\{j,k\},\max\{j,k\}}=mb_j.
\end{equation}
We also make $m$ divisible by all adjunction indices.  The tensor
\begin{equation}
 \label{floor:weighted-product}
 \beta=\bigotimes_{j<k}\beta_{jk}^{\otimes l_{jk}}
\end{equation}
is a nonzero meromorphic $m$-canonical tensor, invariant under
the permutations of the ordered marks.  Evenness removes the
signs from reversed pairs.  It therefore descends to $W$.
More explicitly, relative to the pullback of a local canonical
power frame, its coefficient is the same on all generic sheets.
Factoring $h$ through its finite Stein part, normalized trace
descends this meromorphic coefficient, and modifications do not
change meromorphic functions.  Formula
\eqref{floor:incidence-weights} shows that the descended tensor
has exactly the allowed horizontal poles, of orders $mb_j$.

We must also control \emph{every} vertical prime, including a
prime whose image on $P$ is a point.  Let $E$ be a vertical prime
of $W$, and choose a prime $E^+$ above it, with ramification index
$e$ under $h$.  At its general point the tangential map is
generically finite and separable, so the canonical Jacobian has
order $e-1$.  Near a general point of the image in $P$, take
coordinates $t,s$ with $t$ vanishing on that image.  Whether the
image is a curve or a point,
\[
 \ord_{E^+}(h^*g^*t)\ge e.
\]
In logarithmic differential coordinates at $E^+$, the
differential of this function retains that order.  Both
$d\log x_{jk}$ and the pullback of $ds$ are at most logarithmic.
A logarithmic top form has at most one simple pole.  Writing
$\eta$ in the $dt\wedge ds$ frame, with its holomorphic
coefficient, gives
\begin{equation}
 \label{floor:vertical-order}
 \ord_{E^+}(\beta_{jk})\ge e-1.
\end{equation}
Thus \eqref{floor:weighted-product} has order at least $m(e-1)$,
exactly paying the canonical ramification in descending an
$m$-canonical tensor.  The descended tensor on $W$ is regular at
$E$.

Push it to $Y$.  The horizontal pole orders are allowed by the
effective boundary, and at vertical primes it is regular; hence
normality gives a nonzero section of $mJ_Y$.  Since $J_Y$ is
torsion, this section has no zeros: after trivializing a torsion
power, a nonzero section is a nonzero holomorphic function on a
connected compact normal space.  A vertical floor prime would
give a positive zero in the adjoint section, since its coefficient
is one and the canonical tensor is regular there.  This is
impossible.  Horizontality follows also for $r=0$.

A horizontal conductor surface is generically finite over the
nonprojective surface $P$, so is non-Moishezon.  A Moishezon
neighbor would be projective, and its normalized conductor
surfaces would be projective.  Thus every vertex reached across
a dominant edge remains non-Moishezon.  Finally, each dominant
branch has coefficient one and contributes its generic degree
to the horizontal degree sum two.  This proves the last claim.
\end{proof}

At a non-Moishezon node, if there are two degree-one branches or
one degree-two branch, we consequently have the internal arrows
from \eqref{floor:dlog-residues}, pairing the two marks over $P$.
They are over $Z$ and preserve the restrictions of base sections.
One degree-one branch requires no internal arrow.

\begin{lemma}[Finite loop action for a surface with a curve system]
\label{floor:nonprojective-loops}
Let $S$ be a normal non-Moishezon conductor surface in a cluster
with $r=1$.  A group of $B$-bimeromorphic loops on its semiample
adjunction pair acts through a finite group on every sufficiently
divisible adjoint section space.
\end{lemma}

\begin{proof}
Let $C$ be the normal connected-fiber system base of $J_S$.
It is a polarized projective curve.  A sufficiently divisible
system embeds $C$, and the loop action induces automorphisms of
$C$ preserving a fixed ample power.  Take the smooth minimal
compact K\"ahler surface $Q$ in the bimeromorphic class of $S$.
Classical compact-surface theory will be used in its K\"ahler
and elliptic-surface forms \cite{BHPV}.  Here $a(Q)=1$ and
$H^0(Q,K_Q)\ne0$; bimeromorphic self-maps of this minimal
surface of nonnegative Kodaira dimension are automorphisms.
The map to $C$ is holomorphic on $Q$.  Indeed, on a common
resolved model no curve can dominate $C$ by the nonpositive
N\'eron--Severi test, so the blowdowns to $Q$ factor the map.
Its general fibers are connected.  Every meromorphic function
on $Q$ comes from $C$: the function field has transcendence
degree one and is algebraic over that of $C$, so the function
is constant on connected general smooth fibers and descends
meromorphically.

Let $F$ be the general fiber class.  It is nonzero, rational
and isotropic.  Since the intersection form on $\NS(Q)_\R$
is nonpositive, $K_Q\cdot F=0$.  Adjunction gives genus one
for a general smooth fiber.  The induced cohomology group
preserves the rational line $\Q F$.  On
\begin{equation}
 \label{floor:isotropic-quotient}
 F^\perp/\Q F
\end{equation}
it preserves an integral lattice and the Hodge summands.  The
form is positive definite on real $(2,0)+(0,2)$ and negative
definite on the remaining $(1,1)$ quotient by the Hodge index
theorem.  Reversing the latter sign gives an invariant positive
definite norm.  Its integral isometry group is finite.  In
particular the action on $H^{2,0}(Q)$ is finite.

The action on the base is also finite.  For genus at least two
this follows from finiteness of the curve automorphism group;
for genus one the subgroup preserving a fixed ample line
bundle is finite.  Suppose $C=\PP^1$.  We claim $q(Q)=0$.
Restriction of holomorphic one-forms to a general smooth
elliptic fiber is injective.  Indeed, a form vanishing on one
general fiber vanishes on all general fibers by constancy of
periods, since holomorphic forms on a compact K\"ahler space
are closed.  It is then pulled back from a one-form on the
smooth base open.  This base form extends at a critical value:
at a general point of a fiber divisor the local parameter is
$t=w^d$, and a pole or a nonremovable singularity cannot
pull back to a regular form.  There is no nonzero holomorphic
one-form on $\PP^1$.  Thus $q(Q)\le1$.  If equality held,
the elliptic Albanese map would be nonconstant on a general
elliptic fiber.  A general Albanese fiber would then contain
a curve dominating $\PP^1$, contradicting the absence of
multisections.  This proves the claim.

It follows that
\begin{equation}
 \label{floor:euler-lower}
 \chi(\OO_Q)=1+h^{2,0}(Q)\ge2,
 \qquad e(Q)=12\chi(\OO_Q)-K_Q^2\ge24.
\end{equation}
There are at least three critical values of the relatively
minimal elliptic fibration.  Otherwise the smooth base is
$\PP^1$ with at most two points removed.  Its universal cover
is compact or parabolic, so the genus-one period map to the
upper half-plane is constant.  The integral linear monodromy
is then finite.  Kodaira's singular-fiber and monodromy
classification, including multiple fibers, bounds the Euler
number of a fiber with finite linear monodromy by ten.
The types $I_n$ and their multiples for $n>0$, and $I_n^*$
for $n>0$, have infinite monodromy.  Euler addition with at
most two critical values would therefore give $e(Q)\le20$,
contrary to \eqref{floor:euler-lower}.  The finite set of
critical values is preserved by the base action, and an
automorphism of $\PP^1$ fixing three points is the identity.
The base image is finite in this final case too.

Pass to the finite-index kernel fixing both $C$ and all
holomorphic top forms.  Choose a nonzero holomorphic top form
$\eta$ on $Q$.  Any meromorphic $m$-canonical tensor, divided
by $\eta^{\otimes m}$, is a meromorphic function and hence
comes from $C$.  The kernel fixes it.  It therefore fixes all
the adjoint section spaces under their meromorphic canonical
realizations.  The original loop action has finite image.
\end{proof}

\begin{lemma}[Matching in non-Moishezon clusters with $r=1$]
\label{floor:nonprojective-r1}
Such a cluster admits the required nonzero matching tuple,
with zero restriction on every nondominant branch.
\end{lemma}

\begin{proof}
Use the dominant surfaces as objects, with conductor matchings
and the internal arrows just constructed.  Lemma
\ref{floor:nonprojective-loops} supplies loop finiteness, so
Lemma~\ref{floor:products} gives invariant nonzero sections
in a common degree $m'$.

We verify their descent at each node.  A degree-one branch
$S$ is bimeromorphic to $P$, so its general fibers over $Z$
are connected.  Its section of $m'J_S$ is therefore pulled
back from $m'H$ on $Z$.  When there are two such branches,
the internal arrow makes these base sections identical.
For a degree-two branch, resolve $S\dashrightarrow P$.
Over a general $z\in Z$, every component of the fiber
dominates the connected smooth curve $P_z$, with total
degree two.  Indeed, an exceptional or bad curve on $P$
cannot dominate $Z$, by the no-multisection observation.
Divide the section by the pullback of a local frame of
$m'H$.  This holomorphic function is constant on each
compact connected component of the fiber.  Invariance under
the sheet-exchange arrow equates the two generic values.
It descends meromorphically to $Z$, and then holomorphically:
a meromorphic function on a normal base whose pullback by a
proper surjection is regular has no pole, as can also be
checked through finite Stein factorization.

Pull these base sections back to $Y$.  Their restrictions
match along dominant edges because the object sections do.
They vanish on the prescribed nondominant images because
the invariant object sections retain the initial vanishing
factor and dominate $Z$.  This gives the desired tuple.
\end{proof}

\begin{lemma}[Matching in non-Moishezon clusters with $r=0$]
\label{floor:nonprojective-r0}
Such a cluster also admits a nonzero matching tuple.
\end{lemma}

\begin{proof}
All node adjoints are torsion.  Choose a common sufficiently
divisible even degree and a trivializing section at each
node.  Conductor comparison across an edge is multiplication
by a nonzero constant on these sections.  A graph with no
cycles permits all comparisons to be solved by rescaling.
In general it is enough to prove that the oriented product
of these constants around each cycle is a root of unity,
and then to take one further common power.

The graph has degree at most two by
Lemma~\ref{floor:horizontality}.  At a node on a cycle,
including a cycle formed by parallel edges, the two distinct
branches must each have degree one over its $P_i$.  The
section from the proof of Lemma~\ref{floor:horizontality}
is now a power of $d\log x\wedge\eta_i$: the two unit
points exhaust the horizontal weighted degree.  Its even
residues on the two branches are the corresponding powers
of the base top form $\eta_i$.  Following the cycle gives
bimeromorphic identifications of the $P_i$ and hence
isomorphisms of their smooth minimal models $Q_i$, since
they have nonnegative Kodaira dimension.  The cycle
constant is the scalar on $\eta^{\otimes m}$, or its
inverse, for the resulting automorphism $\phi$ of one
minimal model $Q$.

The existence of a single ambient K\"ahler class restricts
this automorphism.  Fix a K\"ahler form $\omega$ on $V$.
At the two branches of a cycle node, choose smooth branch
resolutions and maps
\[
 r_k:B_k\longrightarrow W_i,\qquad
 \alpha_k:B_k\longrightarrow Q_i,\qquad k=1,2,
\]
where $\alpha_k$ is bimeromorphic.  Put
\[
 v_k=(\alpha_k)_*r_k^*p_i^*[\omega]
       \in H^{1,1}(Q_i,\R).
\]
Each $v_k$ has positive square.  In fact, the pulled-back
ambient class $\beta_k$ on $B_k$ has positive square because
the branch maps generically birationally onto its surface
image in $V$.  Write
$\beta_k=\alpha_k^*v_k+e_k$ with $e_k$ exceptional.
Orthogonality and negative definiteness of exceptional
surface classes give
\begin{equation}
 \label{floor:positive-branch-class}
 v_k^2=\beta_k^2-e_k^2>0.
\end{equation}

Furthermore,
\begin{equation}
 \label{floor:node-congruence}
 v_1-v_2\in\NS(Q_i)_\R.
\end{equation}
Indeed, the rational Hodge morphism
\[
 (\alpha_1)_*r_1^*-(\alpha_2)_*r_2^*
 :H^2(W_i,\Q)\longrightarrow H^2(Q_i,\Q)
\]
kills $H^{2,0}$ and $H^{0,2}$.  Holomorphic two-forms are
basic on the smooth rational-fiber open, by the differential
sequence and the absence of fiberwise one-forms; their
restrictions to the two birational sections agree.  Hence
the rational image of this Hodge morphism has type $(1,1)$,
and is contained in the rational N\'eron--Severi space.
Apply it to the ambient real class to obtain
\eqref{floor:node-congruence}.  Across a conductor edge the
corresponding branch maps to $V$ agree on a common higher
model.  The transported $v_k$ therefore agree across that
edge.  Going around the cycle yields
\begin{equation}
 \label{floor:cycle-congruence}
 \phi^*v-v\in N:=\NS(Q)_\R,\qquad v^2>0.
\end{equation}

The form on $N$ is nonpositive.  If it is negative definite
(including $N=0$), project $v$ orthogonally to $N^\perp$.
The result $w$ has positive square and is fixed by $\phi^*$,
by \eqref{floor:cycle-congruence}.  The Hodge index theorem
makes the $(1,1)$ complement of $w$ negative definite;
the real $(2,0)+(0,2)$ summand is positive definite.
Together these give a positive definite real norm on
$H^2(Q,\R)$ preserved by the integral action.  The cyclic
image is finite, so in particular its eigenvalues on
$H^{2,0}(Q)$ are roots of unity.  If $N$ has a radical,
the Hodge index theorem makes it a rational invariant
isotropic line.  Its orthogonal quotient has the definite
Hodge summands and integral lattice used in
\eqref{floor:isotropic-quotient}; the image on $H^{2,0}$
is again finite.

These exhaust the possibilities for $N$.  Since
$\phi^*\eta^{\otimes m}$ is proportional to
$\eta^{\otimes m}$, the nonzero form $\eta$ itself is an
eigenform: its meromorphic ratio to its pullback has an
$m$th power which is constant, and hence is constant.
Its eigenvalue, and thus the cycle constant, is a root of
unity.  Taking a common power over the finitely many
cycles and then rescaling the node sections solves all
conductor comparisons.
\end{proof}

\subsection{Projective clusters and log-trivial linking}

Every remaining cluster consists of projective vertices.  Indeed,
Lemma~\ref{floor:horizontality} excludes a dominant edge from a
non-Moishezon vertex to a Moishezon one, and a Moishezon vertex is
projective by the klt-type and Namikawa argument above.  All the
strata and birational diagrams used in such a cluster are
algebraic in characteristic zero.

For loop finiteness we use the projective
$B$-pluricanonical-representation theorem: for a projective lc
pair with effective rational boundary and semiample rational
adjoint, its $B$-birational group acts with finite image on
sufficiently divisible log-pluricanonical systems
\cite[Theorem~1.1]{FujinoGongyoPluricanonical}.  This applies to
our normal projective stratum objects with their full
adjunction boundaries.  The remaining issue is to specify enough
internal arrows so that the invariant object sections descend
to the nodes.  We give the required linking arguments.

\begin{lemma}[Connected components of the non-klt locus]
\label{floor:nklt-components}
Let two effective projective lc pairs be crepant bimeromorphic,
and suppose each is klt after decreasing its floor.  On a common
projective log resolution write their common crepant
subboundary as $\Theta$, and put
\[
 T_\Theta=\Theta^{=1},\qquad
 N=-\lfloor\Theta^{<1}\rfloor.
\]
Then $N$ is effective and exceptional for each projection, and
$T_\Theta$ maps with connected fibers onto the non-klt locus
of either pair.  In particular their floor supports have the
same number of connected components.
\end{lemma}

\begin{proof}
Fix either projection $d$.  Nonexceptional coefficients are
those of an effective subboundary, so the negative part
contributing to $N$ is exceptional.  The divisor
\[
 N-T_\Theta=-\lfloor\Theta\rfloor
\]
differs from the smooth klt adjoint $K+\{\Theta\}$ by
$-(K+\Theta)$, a pullback from the base.  It is thus relatively
nef and big; relative bigness here uses that $d$ is birational.
Relative Kawamata--Viehweg vanishing gives
\begin{equation}
 \label{floor:kv-connectedness}
 R^1d_*\OO(N-T_\Theta)=0
\end{equation}
\cite{KawamataViehwegVanishing}.  Normality gives
$d_*\OO(N)=\OO$.  The divisor sequence and
\eqref{floor:kv-connectedness} therefore give a surjection
\[
 \OO\longrightarrow d_*\OO_{T_\Theta}(N).
\]
It factors through $d_*\OO_{T_\Theta}$, which injects into
$d_*\OO_{T_\Theta}(N)$ because $N$ has no component in
$T_\Theta$.  Thus $d_*\OO_{T_\Theta}$ is the structure
sheaf of the reduced image of $T_\Theta$.  This image is
the non-klt locus.  Stein factorization, or the absence of
nontrivial fiberwise idempotents in this equality, shows
that the fibers are connected.  A proper surjection with
connected fibers induces a bijection on connected
components.  The common $T_\Theta$ gives the assertion
for both projections.  Since the pairs become klt after
decreasing the floor, their non-klt loci are exactly their
floor supports.
\end{proof}

We also use the following elementary consequence of relative
degree zero.  Let $E\ge0$ be a rationally Cartier divisor on a
connected projective fiber of a contraction, with $E\cdot C=0$
for every contracted curve $C$.  If the fiber meets $E$, its
whole support is contained in $E$.  Otherwise a chain of
projective curves joining the part in $E$ to a point outside
it contains a curve not in $E$ which meets $E$.  Its
intersection with $E$ is positive, a contradiction.  We shall
refer to this as the \emph{whole-fiber property}.  In particular,
a vertical floor prime of degree zero which meets a fiber
meets every horizontal floor component, since each horizontal
component surjects onto the contraction base.

\begin{lemma}[Log-trivial internal linking in dimension at most three]
\label{floor:trivial-linking}
Let $(Y,\Gamma)$ be a connected normal projective dlt pair of
dimension at most three, with effective rational boundary,
nonempty floor, and
$K_Y+\Gamma\sim_\Q0$.  Suppose the underlying space admits
an effective rational klt pair.  Then the normalized minimal
lc centers, with chain adjunction, can be connected by
$B$-birational comparisons preserving the canonical residues
of a common trivializing log-pluricanonical section in
sufficiently divisible even degrees.
\end{lemma}

\begin{proof}
Take a projective small $\Q$-factorialization for the
underlying klt pair \cite{BCHM}.  It is crepant for the full
adjunction pair and isomorphic at general SNC stratum
points.  The full pair remains dlt, its centers correspond
bimeromorphically, and the chain adjunctions and even residue
maps correspond on common resolutions.  We may therefore
work on this $\Q$-factorial model.

Proceed by dimension.  In dimension one, a unit marking on
a log-trivial curve forces a rational curve; there are at
most two unit markings.  For two markings the even residues
of the logarithmic generator agree by
\eqref{floor:dlog-residues}.  A single marking requires no
comparison.

Suppose first that $\lfloor\Gamma\rfloor$ is connected.
Each normalized floor component with full adjunction again
has a dlt log-trivial pair, and it admits a klt pair by
adjunction keeping only that component in the ambient
boundary.  Within it use induction if there are proper
lc subcenters; otherwise that component is itself the
minimal lc center.  For two successive intersecting floor
components choose a minimal lc center in their intersection.
Such centers exist by ordinary dlt stratum adjunction;
intersecting distinct unit primes meet in lc strata, as is
also seen on a general transverse surface.  The two even
chain restrictions through this shared center agree.
Connectedness of the floor then links all the minimal
centers.

Suppose instead that the floor is disconnected.  Decrease
all its coefficients by one fixed small positive rational
$\epsilon$.  The resulting pair is klt, and its adjoint
is $\Q$-linearly equivalent to
$-\epsilon\lfloor\Gamma\rfloor$, so is not
pseudo-effective.  The projective klt MMP in dimension at
most three ends in a $\Q$-factorial elementary Mori fiber
space
\[
 Y'\longrightarrow P
\]
\cite{ThreefoldFlipsAbundance,BCHM}.  These steps are
crepant for the full log-trivial pair.  One may see this
by transporting a trivializing pluricanonical tensor:
there is no extraction, its divisor remains the negative
of the corresponding multiple of the pushed full boundary,
and the resulting pullback formulas agree.  Equivalently,
the exceptional comparison of the full, relatively
numerically trivial adjoints is zero by negativity.

The pushed floor is relatively ample and supports the
entire non-klt locus, because its decrease is klt.
Lemma~\ref{floor:nklt-components} shows that it is still
disconnected.  Some floor component is horizontal by
relative ampleness.  A vertical floor prime has degree
zero on the Mori ray, since it misses a general fiber.
By the whole-fiber property it contains every fiber it
meets, and so meets every horizontal floor component.
Every other vertical prime is likewise connected to the
horizontal components.  The presence of any vertical
floor would consequently connect the entire floor.
There are therefore no vertical floor primes.

Disconnectedness is now visible on a general fiber, since
every component is horizontal.  That fiber must be a
curve.  Indeed, on a normal projective Cohen--Macaulay
fiber of dimension at least two an effective ample
divisor has connected support.  To check this, take a
large Cartier multiple $A$ of it.  Serre duality and
Serre vanishing give $H^1(\OO(-kA))=0$ for $k\gg0$;
the divisor sequence then gives
$H^0(\OO_{kA})=\C$, implying connected support.
The general fiber here is of klt type and hence
Cohen--Macaulay, and the floor restricts to an effective
ample divisor.  Its restriction is lc and becomes klt
after decreasing the floor, as can be checked on the
restricted common log resolution.

The general fiber is consequently a rational curve with
exactly two degree-one unit markings: it has a nonempty
boundary and log-trivial adjoint, the total weighted
degree is two, and disconnectedness gives two different
horizontal components.  Over a smooth projective model
$R$ of $P$ we obtain a bimeromorphic model
$\PP^1\times R$ with these marks as zero and infinity.
The trivializing tensor is, over the function field and
hence meromorphically, a power of the two-pointed
$d\log$ generator times a meromorphic pluricanonical
tensor on $R$.  Resolving its divisor on $R$, the
crepant subboundary on the product is therefore
\begin{equation}
 \label{floor:product-subboundary}
 \{0\}\times R+\{\infty\}\times R
       +\operatorname{pr}_R^*B_R,
\end{equation}
where $B_R$ may be signed, is log smooth, and has all
coefficients at most one.

No vertical coefficient in
\eqref{floor:product-subboundary} can equal one.  Such
a prime would connect the two sections by an SNC path
of unit components.  This path remains connected
through unit components on a common resolution mapping
to $Y'$.  Explicitly, when a smooth blowup separates
two adjacent unit components at their general
codimension-two crossing, its exceptional divisor has
coefficient one and joins their strict transforms.
Blowups not separating that crossing leave the
connection.  The image of the path would lie in the
non-klt locus of $Y'$ and meet both disconnected floor
components, a contradiction.

All vertical coefficients are therefore strictly below
one.  The log-smooth discrepancy formula now says that
the only zero-log-discrepancy places are the two
horizontal sections themselves.  In fact zero places
of a log-smooth subboundary with coefficients at most
one are generated by its unit strata; here the two
unit divisors are disjoint and there are no deeper
unit strata.  Both section valuations occur as
divisors on the original pair, since its floor is
disconnected.  They are exactly its minimal lc
centers.  Pair them bimeromorphically through $R$.
Their residue restrictions from the common generator
agree in even powers by \eqref{floor:dlog-residues}.
Since these restrictions trivialize their adjunction
lines, equality as meromorphic tensors on a common
resolution gives $B$-crepancy.  This finishes the
disconnected case and the induction.
\end{proof}

\begin{lemma}[Matching in projective clusters with $r=0$]
\label{floor:projective-r0}
Every projective cluster with $r=0$ has a nonzero matching
tuple.
\end{lemma}

\begin{proof}
Use as objects the normalized minimal lc centers of the
nodes, with their iterated full adjunction pairs.
Their adjoints are torsion.  Lemma
\ref{floor:trivial-linking} supplies internal arrows
connecting all objects at a node and preserving even
restrictions from a trivializing node section.

For a matched conductor surface choose a minimal lc
center inside it.  More precisely, choose a minimal
ambient lc stratum contained in the double stratum.
It is a minimal center for the chains on both sides:
any further subcenter would again be a nested ambient
stratum, by the generic SNC description.  Its shared
normalization gives an arrow between the two node
objects, with the symmetric even residue comparison
from $V$.

The projective representation theorem gives finite
loop images.  Lemma~\ref{floor:products} therefore
provides nonzero invariant sections in a common
degree.  They are trivializing sections on the
objects.  At a fixed node, their scalar ratios to
restrictions of one trivializing node section agree
by the internal arrows, so one node section induces
all of them.  Across a conductor surface the two
node restrictions are trivializations of the same
torsion line.  Their ratio is constant on that
connected normal surface.  It is one on the chosen
minimal center by the even residue comparison, so
it is one everywhere.  The node sections thus match
on every conductor surface.
\end{proof}

The preceding linking argument also explains the
usual construction of pre-admissible sections; compare
\cite[Section~5]{GongyoNumericallyTrivial}
and \cite{KollarSources}.  We have given it here to retain
the actual residue comparisons needed for the present
ambient line.

\begin{lemma}[Matching in projective clusters with $r=2$]
\label{floor:projective-r2}
Every projective cluster with $r=2$ admits a nonzero
matching tuple vanishing on all nondominant branches.
\end{lemma}

\begin{proof}
At a node with a dominant branch, the general
$f_i$-fiber is a smooth connected rational curve.
Indeed it has
$K_F+\Gamma_F\sim_\Q0$ and a unit marking, so
adjunction forces genus zero and horizontal weighted
boundary degree two.  Use dominant surface branches
as objects.  If there are two degree-one unit
branches, or one degree-two unit branch, the
two-mark calculation
\eqref{floor:dlog-residues} gives their internal
comparison or sheet-exchange involution over $Z_i$.
These are $B$-crepant and preserve restrictions of
base sections.  A single degree-one unit branch
needs no internal comparison, even if other
horizontal markings have fractional coefficients.

Combine these arrows with conductor matchings and
apply the projective representation theorem and
Lemma~\ref{floor:products}.  At a degree-one object,
a section descends to the base by the birational
comparison.  At a degree-two object, divide by a
local base frame.  Invariance under the involution
equates the two generic values, so the ratio is a
meromorphic base function; normality and properness
give holomorphic descent.  For two degree-one
objects the internal arrow equates their descended
sections.  Pulling back gives the desired node
sections.  Conductor matchings and the retained
initial vanishing factors give exactly the required
matching and nondominant vanishing.
\end{proof}

\subsection{Projective clusters with a curve system}

It remains to treat $r=1$.  Fix a projective node
$(Y,\Gamma)$ with system map $f:Y\to Z$, where $Z$
is a projective curve.  On sufficiently general fibers,
including fibers on simultaneous log resolutions, we
have a connected normal projective dlt log surface
\begin{equation}
 \label{floor:surface-fiber}
 (F,\Gamma_F),\qquad
 K_F+\Gamma_F\sim_\Q0,\qquad
 \Sigma_F:=\lfloor\Gamma_F\rfloor\ne0.
\end{equation}
Normality follows from the general-slice and
Cohen--Macaulay properties, and dlt follows from the
restricted discrepancy formula and general
transversality.  Horizontal boundary components
restrict reduced, while vertical components are
absent on this general fiber.  The floor curves are
smooth and their intersections are ordinary double
lc points.  Effective curve adjunction shows that a
floor curve meeting another is rational and meets
the other floor curves at at most two points.  If
there are two points, its entire log boundary is
precisely these two unit markings.

\begin{lemma}[Disconnected floor on a general surface fiber]
\label{floor:disconnected-fiber}
If $\Sigma_F$ in \eqref{floor:surface-fiber} is
disconnected, it consists of exactly two disjoint
floor curves and there is no deeper lc center in
$F$.  The corresponding dominant floor data on
$Y$ admit an internal pairing over an intermediate
projective surface $P\to Z$ with connected general
fibers.  Generically they are the two unit markings
of log rational-curve fibers on a model crepant
over $Z$.  The two curves over a general $z\in Z$
are bimeromorphic to the same smooth connected
curve $P_z$.  Globally the pairing is either
between two surfaces or an involution of one
surface.
\end{lemma}

\begin{proof}
Start on a projective crepant $\Q$-factorial dlt
model of $Y$.  Decrease all floor coefficients
by a fixed small rational $\epsilon>0$ and run
the projective klt threefold MMP over $Z$.
The decreased adjoint restricts to a negative
multiple of the nonzero effective floor on
a general fiber, so it is not relatively
pseudo-effective.  The program ends in a
relatively elementary Mori fiber contraction
\begin{equation}
 \label{floor:relative-mori}
 Y'\longrightarrow P\longrightarrow Z,
\end{equation}
with $Y'$ $\Q$-factorial.  The full adjoint
remains the actual pullback of $H$ at every
step: all steps are over $Z$ and extract no
divisors, and negativity kills the exceptional
relatively numerically trivial comparison.
Thus the full lc data are crepant throughout.
The transformed floor $\Sigma'$ is relatively
ample over $P$ and supports the full non-klt
locus, since its decrease is klt.  Both maps
in \eqref{floor:relative-mori} have connected
fibers; for the second this also follows by
pushing the structure sheaf of the connected
system map through the first.

On a common smooth projective resolution of
the general fibers $F$ and $F'=Y'_z$, the
crepant subboundary is the same.  Lemma
\ref{floor:nklt-components}, applied to these
surfaces, shows that the floor support on
$F'$ remains disconnected.  This application
does not require that the transformed full
pair be dlt: it is lc, and its decrease is
klt, which are the hypotheses used in that
lemma.

If $\dim P=1$, the floor on $F'$ would be
ample and connected.  One may use the
connected-support argument in the proof of
Lemma~\ref{floor:trivial-linking}, or the
surface Hodge index theorem: a disjoint
splitting of an ample rational divisor would
give orthogonal classes with positive square.
Consequently $\dim P=2$.  A general Mori
fiber over $P$ is rational, with total
horizontal weighted boundary degree two and
at least one unit marking.  A vertical floor
prime has degree zero on the Mori ray.  If
it dominates $Z$, the whole-fiber property
gives, on a general $F'$, a full fiber over
a point of $P_z$, meeting every horizontal
floor curve.  Every other vertical floor
contribution likewise joins the horizontal
ones.  All components restricted from a
horizontal surface dominate $P_z$ for
general $z$, after discarding the finitely
many nongeneral base values.  Thus such a
vertical floor prime would connect the
entire floor of $F'$, which is impossible.

There are therefore no vertical floor
contributions on general $F'$.  Disconnectedness
and the degree sum two force exactly two
disjoint degree-one floor curves over $P_z$.
The curve $P_z$ is smooth and connected for
general $z$: the normal surface base has
only isolated singularities, and generic
smoothness and connected fibers apply.

We also need to exclude hidden deeper
zero-discrepancy places.  Choose the rational
fiber coordinate $x$ over $P_z$ with these
two curves as zero and infinity.  This gives
a bimeromorphic product model
$\PP^1\times P_z$ of $F'$.  A divisible
pluricanonical tensor trivializing its full
log adjoint is, in relative notation,
\begin{equation}
 \label{floor:surface-dlog-product}
 (d\log x)^{\otimes m}\otimes\gamma,
\end{equation}
where $\gamma$ is a meromorphic
$m$-canonical tensor on $P_z$.  There are
no other horizontal boundary terms because
the two unit marks exhaust the degree.
The crepant subboundary on the product is
the two unit sections and vertical fibers
with coefficients
$-\ord_p(\gamma)/m\le1$.

A coefficient-one vertical fiber would
join the two sections by an SNC unit path.
Successive point blowups preserve this
path through unit components: a blowup at
a crossing of two unit components inserts
a unit exceptional curve.  On a common
resolution its image in $F'$ would join
the two disconnected floor curves inside
the non-klt locus.  Thus every vertical
coefficient is strictly below one.  The
log-smooth discrepancy calculation now
leaves only the two horizontal sections
as zero-log-discrepancy places.  They
must both appear as divisors on $F$, since
its floor is disconnected.  Hence its
floor consists of those two curves and
has no deeper lc center.

In particular, no zero-discrepancy place
horizontal over $Z$ is extracted or lost
between $Y$ and $Y'$: restricting such a
place to a general surface fiber would
give another zero-discrepancy place there.
The two Mori marks consequently correspond
bimeromorphically to the original dominant
floor data.  Pair them over $P$, using
normalization of the degree-two cover when
the two marks belong globally to one
surface.  The full crepant equalities over
$Z$ transport local base frames, and the
even $d\log$-residue calculation gives
the asserted internal $B$-comparison.
\end{proof}

Classify a node as \emph{type I} if it has a
deeper curve lc center dominating $Z$, and as
\emph{type II} otherwise.  In type I the floor
of the general surface fiber is connected and
has crossings, by
Lemma~\ref{floor:disconnected-fiber}.  Every
dominant floor surface contains a deeper
dominant center: every vertex of the connected
floor graph meets another, and its crossings
trace such centers.  In type II the floor is
either a single smooth curve, or the
disconnected pair of that lemma.  These types
are constant across dominant edges.  A shared
surface has the same full adjunction on both
sides; a further lc center which dominates
the system curve on one side does so on the
other, equivalently because its curve
adjoint, the restriction of $J$, has positive
degree.  The generic SNC nesting identifies
the further center on both chains.  Since
every branch of a type I node contains such
a center, the type propagates through the
cluster.

\begin{lemma}[Type II matching]
\label{floor:projective-r1-II}
A type II projective cluster with $r=1$
admits a nonzero matching tuple with the
required nondominant vanishing.
\end{lemma}

\begin{proof}
Use dominant surface branches as objects.
For a single smooth floor curve on a
general fiber, the corresponding surface
has connected general fibers over $Z$,
so its sections come from $Z$.  In the
disconnected case use the internal arrows
of Lemma~\ref{floor:disconnected-fiber}.
Over a general $z$, the two curves are
bimeromorphic to the connected curve
$P_z$.  The values of a section relative
to a local base frame are constant on
each of them, and internal invariance
equates the constants.  This remains true
when the two curves are in one global
surface and the arrow is an involution.
Thus invariant sections descend
meromorphically and then holomorphically
to $Z$, using normality and properness.

The projective representation theorem
and Lemma~\ref{floor:products} provide
the invariant nonzero object sections.
Their descended base sections pull back
to node sections.  Matching follows
from the conductor arrows, and vanishing
from the retained initial base factors.
\end{proof}

\begin{lemma}[Type I matching by flagged curves]
\label{floor:projective-r1-I}
A type I projective cluster with $r=1$
also admits a nonzero matching tuple with
the required nondominant vanishing.
\end{lemma}

\begin{proof}
The objects now are normalized curve
strata, \emph{flagged} by the node $Y_i$,
a dominant surface branch $S$ at that
node, and a further normalized curve
branch on $S$ dominating $Z_i$.  The
flag records the actual chain of residue
maps.  Give each curve its full effective
adjunction boundary and the actual
restricted line $J_L$.

There are three kinds of arrows.
First, identify the corresponding
flags across a shared conductor surface.
Second, identify the flags through the
two surface branches at a deeper curve
center within one node.  The generic SNC
calculation and even residues identify
these chains.  Third, consider a fixed
$S$ and the curve Stein factor of
$S\to Z_i$.  If a general connected
fiber has two floor crossings, pair
these two markings over that Stein
curve.  They give a bimeromorphic map
between curve objects or an involution
of one object.  By the surface-fiber
description preceding
Lemma~\ref{floor:disconnected-fiber},
this fiber is log rational with exactly
these two unit points.  Thus
\eqref{floor:dlog-residues}, now applied
to the full adjunction on $S$, gives a
$B$-crepant arrow preserving restrictions
of base sections.  A floor fiber curve
with only one crossing requires no
internal pairing of points.

These arrows connect all flagged
crossing points over a general
$z\in Z_i$ within a node.  Indeed, the
whole floor graph of $F$ is connected,
each floor curve meets another at one
or two points, the two crossings on a
curve are paired by the third kind of
arrow, and the two flags at a crossing
are paired by the second.  This connects
the flags along every path in the
connected graph.

All curve objects are projective lc
pairs with semiample adjoints, so the
projective representation theorem gives
finite loop images.  Apply
Lemma~\ref{floor:products}, with initial
sections pulled from the respective
node bases and vanishing on the required
nondominant images.  In the common final
degree, divide the invariant curve
sections by a local pulled-back base
frame.  The preceding connectivity says
that their values coincide on all flag
sheets over a general $z$.  They
therefore define one meromorphic base
section.  It is holomorphic, since its
pullbacks to the finite dominating
curve covers are holomorphic and the
base is normal.  Pulling it back gives
a section on the node.

Across a dominant conductor surface
$S$, equality on a chosen dominant
further curve implies equality of the
two node restrictions on all of $S$.
To see this precisely, both restrictions
are sections of the semiample actual
line $mJ_S$, hence are pulled from its
connected-fiber Stein system base.
That base is a curve, and the chosen
curve stratum dominates it, since it
dominates $Z_i$.  Pullback of sections
to that stratum is injective.  The
canonical even adjunction compares
these pullbacks, which agree by the
first kind of arrow; hence the original
restrictions agree.  Finally, the
initial vanishing factors and dominance
of the objects show that the node
sections vanish on every nondominant
branch.  This proves the claim.
\end{proof}

\subsection{Completion of whole-boundary nonvanishing}

\begin{proof}[Proof of Proposition~\ref{floor:section}]
Take a cluster of vertices of maximal system
dimension $r$.  If it is isolated, choose a
nonzero high base section vanishing on all
nondominant branch images, as explained
above.  Otherwise $r\le2$.  A cluster
containing a non-Moishezon vertex is entirely
non-Moishezon and has $r\le1$; Lemmas
\ref{floor:nonprojective-r1} and
\ref{floor:nonprojective-r0} handle its two
possibilities.  Every other cluster is
projective.  Lemmas
\ref{floor:projective-r0},
\ref{floor:projective-r2},
\ref{floor:projective-r1-II}, and
\ref{floor:projective-r1-I} cover all its
possibilities.

In each case we have, in one sufficiently
divisible even degree, a nonzero tuple of
node sections matching on dominant edges
and vanishing on all nondominant branches.
Put zero on components outside the chosen
cluster and take a further common power
if necessary so that $mJ$ is Cartier on
the ambient $V$.  At every conductor
surface either the two sections match
by construction or both restrictions
are zero.  Lemma~\ref{floor:descent}
therefore descends the tuple to an
actual section of $\OO_D(mJ)$.  It is
nonzero because its restriction to a
node in the chosen cluster is nonzero.
This proves \eqref{floor:goal} on the
whole reduced boundary.
\end{proof}
